\medskip

On considère la suite $\left(u_n\right)$ définie par $u_0 = 3$ et, pour tout entier naturel $n$, par:
\[u_{n+1} = 5u_n - 4n - 3.\]

\smallskip

\begin{enumerate}
\item 
	\begin{enumerate}
		\item Démontrer que $u_1 = 12$.
		\item Déterminer $u_2$ en détaillant le calcul.
		\item À l'aide de la calculatrice, conjecturer le sens de variation ainsi que la limite de la suite $\left(u_n\right)$.
	\end{enumerate}
\item 
	\begin{enumerate}
		\item Démontrer par récurrence que, pour tout entier naturel $n$, on a : \[u_n\geqslant n + 1.\]
		\item En déduire la limite de la suite $\left(u_n\right)$.
	\end{enumerate}
\item On considère la suite $\left(v_n\right)$ définie pour tout entier naturel $n$ par: \[v_n = u_n - n - 1.\]
	\begin{enumerate}
		\item Démontrer que la suite $\left(v_n\right)$ est géométrique.
		
Donner sa raison et son premier terme $v_0$.
		\item En déduire, pour tout entier naturel $n$, l'expression de $v_n$ en fonction de $n$.
		\item En déduire que pour tout entier naturel $n$ :
\[u_n = 2 \times 5^n + n + 1.\]
		\item En déduire le sens de variation de la suite $\left(u_n\right)$.
	\end{enumerate}






















\item 
\begin{minipage}[t]{0.68\linewidth}
On considère la fonction ci-contre, écrite de manière incomplète en langage Python et destinée à renvoyer le plus petit entier naturel $n$
tel que $u_n \geqslant 10^7$.
\begin{enumerate}
		\item Recopier le programme et compléter les deux instructions manquantes.
		\item Quelle est la valeur renvoyée par cette fonction ?
	\end{enumerate}
\end{minipage}	
\hfill
\begin{minipage}[t]{0.27\linewidth}
\begin{tabular}[t]{|l|}\hline
def suite() :\\
\hspace{0.75cm} u = 3\\
\hspace{0.75cm} n = 0 \\
\hspace{0.75cm} while \ldots :\hspace*{0.5cm}\\
\hspace{1.25cm} \quad u = \ldots \\
\hspace{1.25cm} \quad n = n + 1\\
\hspace{0.75cm} return n\\ \hline
\end{tabular}
\end{minipage}
	
\end{enumerate}

\bigskip
